\section{Kurtosis under compression}\label{sec:compressedk}

\subsection{The kurtosis of an edited record}

The kurtosis of a signal compares its fourth moment with the square of its variance. Removing the
quiet stretches of a record raises the variance of what remains; whether it raises the fourth moment as
much depends on where the fourth moment sits. For a modulated signal the computation is immediate with
the local variance of Section~\ref{sec:modulated}: keeping the stretches with $s$ above a threshold,
the moments of what is kept are the conditional moments of $s$, and by~\eqref{eq:kmod}
\begin{equation}\label{eq:kedited}
  \kappa_{\mathrm{edited}} = 3\,\frac{\tavg{s^2 \mid \text{kept}}}{\tavg{s\mid\text{kept}}^2} .
\end{equation}
For the lognormal envelope of Section~\ref{sec:lognormal} the conditional moments are the truncated
expectations computed there, with $c = 2$ and $c = 1$:
\begin{equation}
  \kappa_{\mathrm{edited}}(q) = 3\,e^{\tau^2}\,\frac{\Phi(u+2\tau)\,\Phi(u)}{\Phi(u+\tau)^2},
  \qquad u = \Phi^{-1}(q),
\end{equation}
which tends to $3e^{\tau^2} = \kappa$ when everything is kept ($q\to1$) and decreases as the test is
shortened (last two columns of Table~12.1): a record with kurtosis 6 edited to a quarter of its length
has kurtosis 4.0. The measured values agree: 6.04 to 4.5 on the example record. The reason is the
one stated above: the quiet time contributed little to the fourth moment but much to the normalisation.

The decrease is a property of smooth modulation. When the kurtosis of a record comes from a few very
large shocks, the editing keeps the shocks, the removed quiet time raises the variance only moderately,
and the kurtosis hardly changes: 6.47 to 6.5 on the thesis reference record. The two records that look
alike by their kurtosis behave differently under compression, which is one more reason to look at the
Lorenz curve rather than at the kurtosis.

\subsection{The kurtosis target of a compressed test}

The methods that prescribe the kurtosis (kFDS, RF, and the kurtosis control methods) use by default
the kurtosis of the target record. Section~\ref{sec:editresults} showed what this does under
compression: a test that must carry the damage of the record in a quarter of the time \emph{and} keep
the record's kurtosis is forced to a heavier tail than the record has, exact at the design exponent
and two to three times too damaging at $b = 9$ on the example record. The edited record shows the
natural target: the kurtosis that the damaging part of the record has once the quiet time is gone.

\begin{result}[Compressed kurtosis target]
For a test of duration $T/\alpha$ built from a record, the kurtosis to prescribe is the kurtosis of the
record edited to $T/\alpha$ (Section~\ref{sec:edit}, without gain), not the kurtosis of the whole
record. The option \code{compress\_kurtosis} (interface: ``compressed kurtosis target''; command line:
\code{--compress-kurtosis}) computes it and uses it as the kurtosis target of kFDS, RF, MLV, VSD and PM.
\end{result}

On the example record this single change brought the $b = 9$ ratio of kFDS from 2.96 to 1.7 and
halved its distance from the record's cycle histogram in the tail; on the thesis record, whose kurtosis
does not drop under editing, it changed nothing, and on seven other channels it improved kFDS in most
cases and worsened it in two where the record's kurtosis was low to begin with. It is an option, not a
default: when the Lorenz curve says that a record is compressible, it is the right target; when the
kurtosis comes from shocks it is neutral.

The same reasoning applies to the response kurtosis spectrum $\kappa_z(\fn)$ that MMVSD prescribes.
The $\kappa_z$ of the edited record is lower than that of the whole record at the oscillators where
quiet time was removed, and it is the $\kappa_z$ that a compressed test should carry. Version 2.2 does
not yet pass it to MMVSD automatically; it can be done by hand by analysing an \code{edit} result (no
gain) as a record and using it as the MMVSD target.
